%%=============================================================================
%% Voorwoord
%%=============================================================================

\chapter*{\IfLanguageName{dutch}{Woord vooraf}{Preface}}%
\label{ch:voorwoord}

Tijdens mijn opleiding Toegepaste Informatica aan HOGENT koos ik bewust voor de keuzerichting Mainframe. Daar leerde ik COBOL via uitgebreide handleidingen en klassieke colleges, terwijl ik moderne talen in de jaren ervoor net via interactieve sandboxes en gamified leerpaden had aangeleerd. Dat contrast intrigeerde me, en het werd nog scherper toen ik hoorde dat 250 miljard regels COBOL nog steeds de wereldeconomie aansturen terwijl de generatie die ze onderhoudt met pensioen gaat. Mijn eigen leerproces gaf me een goed beeld van welke COBOL-bouwstenen voor een beginner het lastigst zijn, en die ervaring werd het vertrekpunt van deze bachelorproef.

Tot slot wil ik enkele mensen bedanken. In de eerste plaats mijn co-promotor, dhr.\ S.\ Willems, wiens grondige feedback en scherpe kritische vragen op elke iteratie van dit werk de kwaliteit ervan een aanzienlijke duw in de rug hebben gegeven. Zijn blik vanuit het werkveld hielp telkens opnieuw om vage formuleringen aan te scherpen, redeneringen te onderbouwen en de scope realistisch te houden. Zonder die voortdurende kritische lezing zou deze bachelorproef niet zijn wat hij vandaag is. Daarnaast wil ik ook de docenten en medestudenten van de Mainframe-keuzerichting bedanken, die mijn COBOL-fundament hebben gelegd waarop dit werk verder bouwt. En mijn familie en vrienden, die het hele traject mee uitzaten en op de juiste momenten zorgden voor de nodige afleiding.

\bigskip
\noindent
Veel leesplezier.

\bigskip
\noindent
Silas De Vuyst \\
Gent, mei 2026
